\documentclass[11pt,a4paper]{article}
\newcommand{\gradelabel}{Grade 10}
\newcommand{\leconhdr}{Calculatrice ou jeu interactif}
% Préambule commun aux plans de cours de robotique (Grades 7 à 12).
% Chaque plan de cours commence par \documentclass[11pt,a4paper]{article}, définit
% \gradelabel et \leconhdr (en-têtes de page), puis % Préambule commun aux plans de cours de robotique (Grades 7 à 12).
% Chaque plan de cours commence par \documentclass[11pt,a4paper]{article}, définit
% \gradelabel et \leconhdr (en-têtes de page), puis % Préambule commun aux plans de cours de robotique (Grades 7 à 12).
% Chaque plan de cours commence par \documentclass[11pt,a4paper]{article}, définit
% \gradelabel et \leconhdr (en-têtes de page), puis \input{../../preambule_robotique.tex}.
\usepackage[T1]{fontenc}
\usepackage[utf8]{inputenc}
\usepackage{lmodern}
\usepackage[french]{babel}
\usepackage[margin=2cm,headheight=14pt]{geometry}
\usepackage[most]{tcolorbox}
\usepackage{xcolor,tabularx,booktabs,array,enumitem,fancyhdr,amssymb,tikz}
\usetikzlibrary{arrows.meta,positioning}

\pagestyle{fancy}
\fancyhf{}
\lhead{\small Robotique -- \gradelabel{} -- Novembre 2026}
\rhead{\small \leconhdr}
\cfoot{\thepage}
\setlist{itemsep=1pt,topsep=2pt,leftmargin=*}
\renewcommand{\arraystretch}{1.35}
\newcolumntype{Y}{>{\raggedright\arraybackslash}X}
\newcommand{\cb}{$\square$\,}
\newcommand{\fl}{$\rightarrow$}
\newcommand{\lig}{\noindent\rule{\linewidth}{0.3pt}\\[5mm]}

\newtcolorbox{etape}[3]{enhanced,breakable,colback=#1!6,colframe=#1!70!black,
  coltitle=white,fonttitle=\bfseries,title={#2\hfill #3},boxrule=0.8pt,arc=2mm,
  before skip=8pt,after skip=8pt}
\newtcolorbox{astuce}[1][Différenciation et besoins particuliers]{breakable,
  colback=yellow!12,colframe=yellow!50!black,fonttitle=\bfseries\small,
  title={#1},boxrule=0.5pt,arc=1mm,before skip=6pt,after skip=6pt}
\newtcolorbox{code}[1][Code]{breakable,colback=black!4,colframe=black!60,
  fonttitle=\bfseries\small,title={#1},boxrule=0.5pt,arc=1mm,fontupper=\ttfamily\small,
  before skip=6pt,after skip=6pt}
% Blocs de code coloré (écrire le code sans accents : comme dans l'éditeur)
\newtcblisting{arduino}[1][Code Arduino]{listing only,breakable,colback=black!4,
  colframe=black!60,fonttitle=\bfseries\small,title={#1},boxrule=0.5pt,arc=1mm,
  before skip=6pt,after skip=6pt,
  listing options={language=C++,basicstyle=\ttfamily\small,columns=fullflexible,
  keepspaces=true,showstringspaces=false,commentstyle=\color{green!40!black}\itshape,
  keywordstyle=\color{blue!70!black}\bfseries}}
\newtcblisting{python}[1][Code Python]{listing only,breakable,colback=black!4,
  colframe=black!60,fonttitle=\bfseries\small,title={#1},boxrule=0.5pt,arc=1mm,
  before skip=6pt,after skip=6pt,
  listing options={language=Python,basicstyle=\ttfamily\small,columns=fullflexible,
  keepspaces=true,showstringspaces=false,commentstyle=\color{green!40!black}\itshape,
  keywordstyle=\color{blue!70!black}\bfseries}}
\newcommand{\entete}[2]{%
  \begin{tcolorbox}[colback=#1!10,colframe=#1!70!black,boxrule=0.8pt,arc=2mm]
  {\Large\bfseries #2}\hfill Nom : \rule{4cm}{0.3pt}\quad Date : \rule{2.2cm}{0.3pt}
  \end{tcolorbox}}

% Bandeau de titre : #1 titre, #2 semaine
\newcommand{\titreblock}[2]{%
\begin{tcolorbox}[colback=teal!10,colframe=teal!70!black,arc=2mm,boxrule=1pt]
\centering
{\LARGE\bfseries #1}\\[2mm]
{\large Robotique -- \gradelabel{} -- Plan de cours de 2 heures}\\[1mm]
Enseignant : Francis Ki \quad|\quad #2
\end{tcolorbox}}

% Tableau du minutage : 6 contenus (accroche, mini-leçon, guidée, autonome, mise en commun, exit ticket)
\newcommand{\minutage}[6]{%
\begin{tabularx}{\linewidth}{>{\bfseries}p{2.2cm}p{2.6cm}Y p{3cm}}
\toprule
Minutes & Étape & Contenu & Organisation\\
\midrule
0 -- 10 & Accroche & #1 & Classe entière\\
10 -- 30 & Mini-leçon & #2 & Classe entière\\
30 -- 60 & Pratique guidée & #3 & Collectif puis binômes\\
60 -- 95 & Pratique autonome & #4 & Individuel sur poste\\
95 -- 110 & Mise en commun & #5 & Classe entière\\
110 -- 120 & Exit ticket & #6 & Individuel sur poste\\
\bottomrule
\end{tabularx}}

% Étayage progressif : #1 aide 1, #2 aide 2, #3 aide 3
\newcommand{\aides}[3]{%
\textbf{Étayage progressif de l'aide.}
\begin{center}
\begin{tabularx}{\linewidth}{>{\bfseries}p{2.4cm}Y p{4.4cm}}
\toprule
Niveau d'aide & Contenu & Qui intervient\\
\midrule
Aide 1 & #1 & Le binôme seul\\
Aide 2 & #2 & Le binôme avec la carte d'aide\\
Aide 3 & #3 & L'enseignant\\
\bottomrule
\end{tabularx}
\end{center}}

% Les trois groupes de la pratique autonome : #1 tâche A, #2 tâche B, #3 tâche C
\newcommand{\groupes}[3]{%
\textbf{Consigne commune\,:} chaque élève ouvre le fichier de son groupe, le renomme avec son prénom et travaille seul sur son poste. L'enseignant passe d'abord auprès du groupe A pendant 10 min.

\medskip
\begin{tabularx}{\linewidth}{>{\bfseries}p{3cm}Y p{2.2cm}}
\toprule
Groupe & Tâche sur poste & Annexe\\
\midrule
A : en dessous du niveau & #1 & Annexe C\\
B : au niveau & #2 & Annexe D\\
C : au-dessus du niveau & #3 & Annexe E\\
\bottomrule
\end{tabularx}}

% Mise en commun : #1 groupe A, #2 groupe B, #3 groupe C, #4 conclusion
\newcommand{\miseencommun}[4]{%
\begin{etape}{teal}{5. Mise en commun}{95 -- 110 min}
\begin{enumerate}
 \item (2 min) Retour au calme, écrans mis en veille.
 \item (9 min) Trois passages courts de 3 min, avec valorisation :
 \begin{itemize}
  \item \textbf{Groupe A} : #1
  \item \textbf{Groupe B} : #2
  \item \textbf{Groupe C} : #3
 \end{itemize}
 \item (4 min) Conclusion : #4
\end{enumerate}
\end{etape}}

% Clôture : #1 question 1, #2 question 2
\newcommand{\cloture}[2]{%
\begin{etape}{gray}{6. Clôture : exit ticket}{110 -- 120 min}
Formulaire en ligne ou papier (annexe G) :
\begin{enumerate}
 \item #1
 \item #2
 \item Auto-évaluation : \emph{Je maîtrise / Je progresse / J'ai besoin d'aide.}
\end{enumerate}
Lire les réponses avant la séance suivante pour ajuster les groupes de besoin.
\end{etape}}

% Grille d'observation : #1 #2 #3 intitulés de colonnes
\newcommand{\grille}[3]{%
\subsection*{Grille d'observation (enseignant)}
\begin{tabularx}{\linewidth}{Y p{2.4cm} p{2.4cm} p{2.4cm} Y}
\toprule
Élève & #1 & #2 & #3 & Groupe suivant\\
\midrule
 & & & & \\ & & & & \\ & & & & \\ & & & & \\ & & & & \\
\bottomrule
\end{tabularx}}

% Signature de l'exit ticket
\newcommand{\autoeval}{Je me sens : \cb Je maîtrise \quad \cb Je progresse \quad \cb J'ai besoin d'aide}
.
\usepackage[T1]{fontenc}
\usepackage[utf8]{inputenc}
\usepackage{lmodern}
\usepackage[french]{babel}
\usepackage[margin=2cm,headheight=14pt]{geometry}
\usepackage[most]{tcolorbox}
\usepackage{xcolor,tabularx,booktabs,array,enumitem,fancyhdr,amssymb,tikz}
\usetikzlibrary{arrows.meta,positioning}

\pagestyle{fancy}
\fancyhf{}
\lhead{\small Robotique -- \gradelabel{} -- Novembre 2026}
\rhead{\small \leconhdr}
\cfoot{\thepage}
\setlist{itemsep=1pt,topsep=2pt,leftmargin=*}
\renewcommand{\arraystretch}{1.35}
\newcolumntype{Y}{>{\raggedright\arraybackslash}X}
\newcommand{\cb}{$\square$\,}
\newcommand{\fl}{$\rightarrow$}
\newcommand{\lig}{\noindent\rule{\linewidth}{0.3pt}\\[5mm]}

\newtcolorbox{etape}[3]{enhanced,breakable,colback=#1!6,colframe=#1!70!black,
  coltitle=white,fonttitle=\bfseries,title={#2\hfill #3},boxrule=0.8pt,arc=2mm,
  before skip=8pt,after skip=8pt}
\newtcolorbox{astuce}[1][Différenciation et besoins particuliers]{breakable,
  colback=yellow!12,colframe=yellow!50!black,fonttitle=\bfseries\small,
  title={#1},boxrule=0.5pt,arc=1mm,before skip=6pt,after skip=6pt}
\newtcolorbox{code}[1][Code]{breakable,colback=black!4,colframe=black!60,
  fonttitle=\bfseries\small,title={#1},boxrule=0.5pt,arc=1mm,fontupper=\ttfamily\small,
  before skip=6pt,after skip=6pt}
% Blocs de code coloré (écrire le code sans accents : comme dans l'éditeur)
\newtcblisting{arduino}[1][Code Arduino]{listing only,breakable,colback=black!4,
  colframe=black!60,fonttitle=\bfseries\small,title={#1},boxrule=0.5pt,arc=1mm,
  before skip=6pt,after skip=6pt,
  listing options={language=C++,basicstyle=\ttfamily\small,columns=fullflexible,
  keepspaces=true,showstringspaces=false,commentstyle=\color{green!40!black}\itshape,
  keywordstyle=\color{blue!70!black}\bfseries}}
\newtcblisting{python}[1][Code Python]{listing only,breakable,colback=black!4,
  colframe=black!60,fonttitle=\bfseries\small,title={#1},boxrule=0.5pt,arc=1mm,
  before skip=6pt,after skip=6pt,
  listing options={language=Python,basicstyle=\ttfamily\small,columns=fullflexible,
  keepspaces=true,showstringspaces=false,commentstyle=\color{green!40!black}\itshape,
  keywordstyle=\color{blue!70!black}\bfseries}}
\newcommand{\entete}[2]{%
  \begin{tcolorbox}[colback=#1!10,colframe=#1!70!black,boxrule=0.8pt,arc=2mm]
  {\Large\bfseries #2}\hfill Nom : \rule{4cm}{0.3pt}\quad Date : \rule{2.2cm}{0.3pt}
  \end{tcolorbox}}

% Bandeau de titre : #1 titre, #2 semaine
\newcommand{\titreblock}[2]{%
\begin{tcolorbox}[colback=teal!10,colframe=teal!70!black,arc=2mm,boxrule=1pt]
\centering
{\LARGE\bfseries #1}\\[2mm]
{\large Robotique -- \gradelabel{} -- Plan de cours de 2 heures}\\[1mm]
Enseignant : Francis Ki \quad|\quad #2
\end{tcolorbox}}

% Tableau du minutage : 6 contenus (accroche, mini-leçon, guidée, autonome, mise en commun, exit ticket)
\newcommand{\minutage}[6]{%
\begin{tabularx}{\linewidth}{>{\bfseries}p{2.2cm}p{2.6cm}Y p{3cm}}
\toprule
Minutes & Étape & Contenu & Organisation\\
\midrule
0 -- 10 & Accroche & #1 & Classe entière\\
10 -- 30 & Mini-leçon & #2 & Classe entière\\
30 -- 60 & Pratique guidée & #3 & Collectif puis binômes\\
60 -- 95 & Pratique autonome & #4 & Individuel sur poste\\
95 -- 110 & Mise en commun & #5 & Classe entière\\
110 -- 120 & Exit ticket & #6 & Individuel sur poste\\
\bottomrule
\end{tabularx}}

% Étayage progressif : #1 aide 1, #2 aide 2, #3 aide 3
\newcommand{\aides}[3]{%
\textbf{Étayage progressif de l'aide.}
\begin{center}
\begin{tabularx}{\linewidth}{>{\bfseries}p{2.4cm}Y p{4.4cm}}
\toprule
Niveau d'aide & Contenu & Qui intervient\\
\midrule
Aide 1 & #1 & Le binôme seul\\
Aide 2 & #2 & Le binôme avec la carte d'aide\\
Aide 3 & #3 & L'enseignant\\
\bottomrule
\end{tabularx}
\end{center}}

% Les trois groupes de la pratique autonome : #1 tâche A, #2 tâche B, #3 tâche C
\newcommand{\groupes}[3]{%
\textbf{Consigne commune\,:} chaque élève ouvre le fichier de son groupe, le renomme avec son prénom et travaille seul sur son poste. L'enseignant passe d'abord auprès du groupe A pendant 10 min.

\medskip
\begin{tabularx}{\linewidth}{>{\bfseries}p{3cm}Y p{2.2cm}}
\toprule
Groupe & Tâche sur poste & Annexe\\
\midrule
A : en dessous du niveau & #1 & Annexe C\\
B : au niveau & #2 & Annexe D\\
C : au-dessus du niveau & #3 & Annexe E\\
\bottomrule
\end{tabularx}}

% Mise en commun : #1 groupe A, #2 groupe B, #3 groupe C, #4 conclusion
\newcommand{\miseencommun}[4]{%
\begin{etape}{teal}{5. Mise en commun}{95 -- 110 min}
\begin{enumerate}
 \item (2 min) Retour au calme, écrans mis en veille.
 \item (9 min) Trois passages courts de 3 min, avec valorisation :
 \begin{itemize}
  \item \textbf{Groupe A} : #1
  \item \textbf{Groupe B} : #2
  \item \textbf{Groupe C} : #3
 \end{itemize}
 \item (4 min) Conclusion : #4
\end{enumerate}
\end{etape}}

% Clôture : #1 question 1, #2 question 2
\newcommand{\cloture}[2]{%
\begin{etape}{gray}{6. Clôture : exit ticket}{110 -- 120 min}
Formulaire en ligne ou papier (annexe G) :
\begin{enumerate}
 \item #1
 \item #2
 \item Auto-évaluation : \emph{Je maîtrise / Je progresse / J'ai besoin d'aide.}
\end{enumerate}
Lire les réponses avant la séance suivante pour ajuster les groupes de besoin.
\end{etape}}

% Grille d'observation : #1 #2 #3 intitulés de colonnes
\newcommand{\grille}[3]{%
\subsection*{Grille d'observation (enseignant)}
\begin{tabularx}{\linewidth}{Y p{2.4cm} p{2.4cm} p{2.4cm} Y}
\toprule
Élève & #1 & #2 & #3 & Groupe suivant\\
\midrule
 & & & & \\ & & & & \\ & & & & \\ & & & & \\ & & & & \\
\bottomrule
\end{tabularx}}

% Signature de l'exit ticket
\newcommand{\autoeval}{Je me sens : \cb Je maîtrise \quad \cb Je progresse \quad \cb J'ai besoin d'aide}
.
\usepackage[T1]{fontenc}
\usepackage[utf8]{inputenc}
\usepackage{lmodern}
\usepackage[french]{babel}
\usepackage[margin=2cm,headheight=14pt]{geometry}
\usepackage[most]{tcolorbox}
\usepackage{xcolor,tabularx,booktabs,array,enumitem,fancyhdr,amssymb,tikz}
\usetikzlibrary{arrows.meta,positioning}

\pagestyle{fancy}
\fancyhf{}
\lhead{\small Robotique -- \gradelabel{} -- Novembre 2026}
\rhead{\small \leconhdr}
\cfoot{\thepage}
\setlist{itemsep=1pt,topsep=2pt,leftmargin=*}
\renewcommand{\arraystretch}{1.35}
\newcolumntype{Y}{>{\raggedright\arraybackslash}X}
\newcommand{\cb}{$\square$\,}
\newcommand{\fl}{$\rightarrow$}
\newcommand{\lig}{\noindent\rule{\linewidth}{0.3pt}\\[5mm]}

\newtcolorbox{etape}[3]{enhanced,breakable,colback=#1!6,colframe=#1!70!black,
  coltitle=white,fonttitle=\bfseries,title={#2\hfill #3},boxrule=0.8pt,arc=2mm,
  before skip=8pt,after skip=8pt}
\newtcolorbox{astuce}[1][Différenciation et besoins particuliers]{breakable,
  colback=yellow!12,colframe=yellow!50!black,fonttitle=\bfseries\small,
  title={#1},boxrule=0.5pt,arc=1mm,before skip=6pt,after skip=6pt}
\newtcolorbox{code}[1][Code]{breakable,colback=black!4,colframe=black!60,
  fonttitle=\bfseries\small,title={#1},boxrule=0.5pt,arc=1mm,fontupper=\ttfamily\small,
  before skip=6pt,after skip=6pt}
% Blocs de code coloré (écrire le code sans accents : comme dans l'éditeur)
\newtcblisting{arduino}[1][Code Arduino]{listing only,breakable,colback=black!4,
  colframe=black!60,fonttitle=\bfseries\small,title={#1},boxrule=0.5pt,arc=1mm,
  before skip=6pt,after skip=6pt,
  listing options={language=C++,basicstyle=\ttfamily\small,columns=fullflexible,
  keepspaces=true,showstringspaces=false,commentstyle=\color{green!40!black}\itshape,
  keywordstyle=\color{blue!70!black}\bfseries}}
\newtcblisting{python}[1][Code Python]{listing only,breakable,colback=black!4,
  colframe=black!60,fonttitle=\bfseries\small,title={#1},boxrule=0.5pt,arc=1mm,
  before skip=6pt,after skip=6pt,
  listing options={language=Python,basicstyle=\ttfamily\small,columns=fullflexible,
  keepspaces=true,showstringspaces=false,commentstyle=\color{green!40!black}\itshape,
  keywordstyle=\color{blue!70!black}\bfseries}}
\newcommand{\entete}[2]{%
  \begin{tcolorbox}[colback=#1!10,colframe=#1!70!black,boxrule=0.8pt,arc=2mm]
  {\Large\bfseries #2}\hfill Nom : \rule{4cm}{0.3pt}\quad Date : \rule{2.2cm}{0.3pt}
  \end{tcolorbox}}

% Bandeau de titre : #1 titre, #2 semaine
\newcommand{\titreblock}[2]{%
\begin{tcolorbox}[colback=teal!10,colframe=teal!70!black,arc=2mm,boxrule=1pt]
\centering
{\LARGE\bfseries #1}\\[2mm]
{\large Robotique -- \gradelabel{} -- Plan de cours de 2 heures}\\[1mm]
Enseignant : Francis Ki \quad|\quad #2
\end{tcolorbox}}

% Tableau du minutage : 6 contenus (accroche, mini-leçon, guidée, autonome, mise en commun, exit ticket)
\newcommand{\minutage}[6]{%
\begin{tabularx}{\linewidth}{>{\bfseries}p{2.2cm}p{2.6cm}Y p{3cm}}
\toprule
Minutes & Étape & Contenu & Organisation\\
\midrule
0 -- 10 & Accroche & #1 & Classe entière\\
10 -- 30 & Mini-leçon & #2 & Classe entière\\
30 -- 60 & Pratique guidée & #3 & Collectif puis binômes\\
60 -- 95 & Pratique autonome & #4 & Individuel sur poste\\
95 -- 110 & Mise en commun & #5 & Classe entière\\
110 -- 120 & Exit ticket & #6 & Individuel sur poste\\
\bottomrule
\end{tabularx}}

% Étayage progressif : #1 aide 1, #2 aide 2, #3 aide 3
\newcommand{\aides}[3]{%
\textbf{Étayage progressif de l'aide.}
\begin{center}
\begin{tabularx}{\linewidth}{>{\bfseries}p{2.4cm}Y p{4.4cm}}
\toprule
Niveau d'aide & Contenu & Qui intervient\\
\midrule
Aide 1 & #1 & Le binôme seul\\
Aide 2 & #2 & Le binôme avec la carte d'aide\\
Aide 3 & #3 & L'enseignant\\
\bottomrule
\end{tabularx}
\end{center}}

% Les trois groupes de la pratique autonome : #1 tâche A, #2 tâche B, #3 tâche C
\newcommand{\groupes}[3]{%
\textbf{Consigne commune\,:} chaque élève ouvre le fichier de son groupe, le renomme avec son prénom et travaille seul sur son poste. L'enseignant passe d'abord auprès du groupe A pendant 10 min.

\medskip
\begin{tabularx}{\linewidth}{>{\bfseries}p{3cm}Y p{2.2cm}}
\toprule
Groupe & Tâche sur poste & Annexe\\
\midrule
A : en dessous du niveau & #1 & Annexe C\\
B : au niveau & #2 & Annexe D\\
C : au-dessus du niveau & #3 & Annexe E\\
\bottomrule
\end{tabularx}}

% Mise en commun : #1 groupe A, #2 groupe B, #3 groupe C, #4 conclusion
\newcommand{\miseencommun}[4]{%
\begin{etape}{teal}{5. Mise en commun}{95 -- 110 min}
\begin{enumerate}
 \item (2 min) Retour au calme, écrans mis en veille.
 \item (9 min) Trois passages courts de 3 min, avec valorisation :
 \begin{itemize}
  \item \textbf{Groupe A} : #1
  \item \textbf{Groupe B} : #2
  \item \textbf{Groupe C} : #3
 \end{itemize}
 \item (4 min) Conclusion : #4
\end{enumerate}
\end{etape}}

% Clôture : #1 question 1, #2 question 2
\newcommand{\cloture}[2]{%
\begin{etape}{gray}{6. Clôture : exit ticket}{110 -- 120 min}
Formulaire en ligne ou papier (annexe G) :
\begin{enumerate}
 \item #1
 \item #2
 \item Auto-évaluation : \emph{Je maîtrise / Je progresse / J'ai besoin d'aide.}
\end{enumerate}
Lire les réponses avant la séance suivante pour ajuster les groupes de besoin.
\end{etape}}

% Grille d'observation : #1 #2 #3 intitulés de colonnes
\newcommand{\grille}[3]{%
\subsection*{Grille d'observation (enseignant)}
\begin{tabularx}{\linewidth}{Y p{2.4cm} p{2.4cm} p{2.4cm} Y}
\toprule
Élève & #1 & #2 & #3 & Groupe suivant\\
\midrule
 & & & & \\ & & & & \\ & & & & \\ & & & & \\ & & & & \\
\bottomrule
\end{tabularx}}

% Signature de l'exit ticket
\newcommand{\autoeval}{Je me sens : \cb Je maîtrise \quad \cb Je progresse \quad \cb J'ai besoin d'aide}


\begin{document}
\titreblock{Calculatrice ou jeu interactif (Python)}{Semaine du 23 au 27 novembre 2026}

%==================================================================
\section{Informations générales}
\begin{tabularx}{\linewidth}{>{\bfseries}p{3.6cm}Y}
\toprule
Titre & Calculatrice ou jeu interactif : un programme complet\\
Classe / effectif & Grade 10, 11 à 20 élèves, niveaux hétérogènes, élèves à besoins particuliers\\
Durée & 2 heures (120 min)\\
Lieu & Salle informatique, un poste par élève\\
Matériel et logiciels &
\begin{itemize}
 \item Un poste par élève avec l'éditeur \textbf{Mu} en mode \emph{Python 3} (Vérifier, Exécuter, Débogueur).
 \item Vidéoprojecteur ou tableau interactif avec Mu projeté.
 \item Un dossier partagé contenant les fichiers de départ \texttt{.py} des trois niveaux et la grille d'évaluation du projet.
 \item Inspiration : unité «~Putting it all together~» du cours TeachComputing (binômes, codage en direct, exemples commentés).
\end{itemize}\\
Pré-requis &
Séances précédentes : variables, types, saisie, conditions, boucles \texttt{for} et \texttt{while}.\\
\bottomrule
\end{tabularx}

\begin{astuce}[Préparation avant le cours (environ 45 min)]
\begin{enumerate}
 \item Préparer un exemple fini de calculatrice et un exemple fini de jeu (à montrer en démonstration).
 \item Préparer les fichiers de départ : \texttt{niveauA.py}, \texttt{niveauB.py}, \texttt{niveauC.py}.
 \item Préparer la grille d'évaluation du projet (annexe B) et le formulaire de sortie (annexe G).
 \item Prévoir la séance de décembre : \og Introduction au HTML \fg{}, qui suivra l'évaluation.
\end{enumerate}
\end{astuce}

%==================================================================
\section{Objectifs d'apprentissage}
\begin{tcolorbox}[colback=gray!8,colframe=gray!60!black,arc=1mm,boxrule=0.5pt]
\textbf{Objectifs du programme}
\begin{itemize}
 \item Concevoir un programme fonctionnel.
 \item Corriger les erreurs détectées.
\end{itemize}
\textbf{À la fin de la séance, l'élève est capable\,:}
\begin{itemize}
 \item de décomposer un programme en étapes (pseudo-code) avant de le coder\,;
 \item d'assembler variables, conditions et boucles dans un programme complet\,;
 \item de distinguer et de corriger une erreur de syntaxe, d'exécution et de logique.
\end{itemize}
\textbf{Compétences numériques\,:} planifier, tester par un jeu d'essais, lire un message d'erreur, enregistrer des versions.
\end{tcolorbox}

%==================================================================
\section{Déroulé minuté (120 min)}
\minutage{Démonstration d'une calculatrice et d'un jeu}{Du pseudo-code au programme ; les trois types d'erreurs}{Pseudo-code collectif, puis première opération en binômes}{Trois niveaux de besoin}{Présentation des programmes}{Trois questions sur poste}

%------------------------------------------------------------------
\begin{etape}{orange}{1. Accroche : deux programmes à découvrir}{0 -- 10 min}
\begin{enumerate}
 \item (5 min) Faire fonctionner une calculatrice (menu, boucle, quitter) et un jeu de devinette. \emph{Quelles variables, conditions et boucles y voyez-vous ?}
 \item (5 min) Les élèves choisissent leur projet : calculatrice ou jeu. Question-problème : \textbf{«~Comment transformer une idée en programme qui fonctionne et qui ne plante pas ?~»}
\end{enumerate}
\end{etape}

%------------------------------------------------------------------
\begin{etape}{blue}{2. Mini-leçon : planifier, tester, déboguer}{10 -- 30 min}
\textbf{Support\,:} annexe A. Codage en direct. Quatre temps de 5 min.

\textbf{Temps 1 : du pseudo-code au code.}
\begin{center}
\begin{tikzpicture}[box/.style={draw,rounded corners,thick,minimum height=1cm,text width=2.7cm,align=center,font=\small\bfseries},>={Stealth[length=3mm]},node distance=0.7cm]
\node[box,fill=orange!20] (a) {1. Idée};
\node[box,fill=blue!15,right=of a] (b) {2. Pseudo-code};
\node[box,fill=green!15,right=of b] (c) {3. Code et test};
\node[box,fill=violet!15,right=of c] (d) {4. Corriger};
\draw[->,thick] (a)--(b); \draw[->,thick] (b)--(c); \draw[->,thick] (c)--(d);
\end{tikzpicture}
\end{center}

\begin{python}[Pseudo-code d'une calculatrice]
# repeter jusqu'a ce que l'utilisateur quitte
#   afficher le menu
#   demander l'operation
#   si quitter : arreter
#   demander deux nombres
#   selon l'operation : calculer
#   afficher le resultat
\end{python}

\textbf{Temps 2 : les trois types d'erreurs.}
\begin{center}
\begin{tabularx}{\linewidth}{>{\bfseries}p{3.2cm}Y Y}
\toprule
Type & Exemple & Comment la trouver\\
\midrule
Syntaxe & Deux-points oublié, parenthèse non fermée & Mu la signale avant ou à l'exécution\\
Exécution & Division par zéro, \texttt{int("abc")} & Le programme s'arrête avec un message (\emph{Traceback})\\
Logique & Mauvaise condition, résultat faux & Aucun message : il faut tester\\
\bottomrule
\end{tabularx}
\end{center}

\textbf{Temps 3 : lire un message d'erreur.} Montrer un \emph{Traceback} : la dernière ligne dit le type d'erreur, l'avant-dernière donne la ligne concernée.

\textbf{Temps 4 : le jeu d'essais.} Avant de dire «~ça marche~», tester des cas normaux, des cas limites (0, négatif) et des saisies absurdes.

\begin{astuce}
Fournir le pseudo-code sous forme d'étiquettes à ranger dans l'ordre. Utiliser les couleurs de l'annexe A pour distinguer les trois erreurs. Autoriser le travail en binôme sur la planification.
\end{astuce}
\end{etape}

%------------------------------------------------------------------
\begin{etape}{green}{3. Pratique guidée}{30 -- 60 min}
\textbf{Fiche\,:} annexe B. Programmation en binômes (pilote et navigateur, rôles échangés à mi-temps).

\textbf{3a. Pseudo-code collectif (12 min).} Écrire ensemble le pseudo-code de la calculatrice, puis coder le menu et la première opération (addition).

\textbf{3b. Première opération en binômes (18 min).} Chaque binôme code, selon son projet, soit l'addition et la soustraction avec le menu, soit la partie «~choisir un nombre secret et le comparer~». Tester avec trois entrées.

\aides{Le pseudo-code du 3a et les modèles de code des séances précédentes}{Questions guides : \emph{Quelle est la prochaine étape du pseudo-code ? Quelle structure (condition ou boucle) faut-il ?}}{L'enseignant écrit la première ligne de code avec le binôme}

\begin{astuce}
Insister : coder une étape, tester, puis ajouter la suivante. Laisser ouverte la liste des modèles de code. En cas de blocage prolongé, passer au pseudo-code sur papier.
\end{astuce}
\end{etape}

%------------------------------------------------------------------
\begin{etape}{violet}{4. Pratique autonome différenciée}{60 -- 95 min}
\groupes{Fichier \texttt{niveauA.py} : calculatrice à deux opérations dont les lignes essentielles sont à compléter, avec le pseudo-code fourni et une banque de mots. Tester trois calculs.}{Programmer une calculatrice à quatre opérations avec menu en boucle, gestion de la division par zéro et option «~Quitter~», ou un jeu de devinette complet. Fournir un jeu d'essais.}{Réaliser un projet complet et structuré : calculatrice avec fonctions (\texttt{def addition(a, b):}, etc.) et historique des calculs dans une liste, ou jeu avec rejouer et meilleur score. Gérer les saisies invalides avec \texttt{try} et \texttt{except}.}

\medskip
\textbf{Plan B si un élève termine tôt\,:} tester le programme d'un camarade en cherchant à le faire planter, et noter les bogues trouvés.
\end{etape}

%------------------------------------------------------------------
\begin{etape}{teal}{5. Mise en commun : présentation des programmes}{95 -- 110 min}
\begin{enumerate}
 \item (2 min) Retour au calme.
 \item (9 min) Trois présentations courtes (3 min), avec valorisation :
 \begin{itemize}
  \item \textbf{Groupe A} : un élève montre ses calculs et explique une ligne. Valoriser la compréhension de son programme.
  \item \textbf{Groupe B} : un élève présente son jeu d'essais et un bogue corrigé. Valoriser la rigueur.
  \item \textbf{Groupe C} : un élève présente sa structure (fonctions, gestion d'erreurs). Valoriser la maîtrise.
 \end{itemize}
 \item (4 min) Conclusion : \textit{un programme se planifie en pseudo-code, se teste avec un jeu d'essais et se corrige en distinguant le type d'erreur.} Annoncer la prochaine séance : page Web avec HTML.
\end{enumerate}
\end{etape}

\cloture{Un programme s'exécute mais affiche un mauvais résultat sans message d'erreur : de quel type d'erreur s'agit-il ?}{Cite deux cas que tu testerais pour vérifier une calculatrice.}

%==================================================================
% ANNEXES
%==================================================================
\newpage
\section*{Annexe A : Support résumé de la mini-leçon}
\entete{blue}{Planifier, tester, déboguer}

\begin{tcolorbox}[colback=yellow!10,colframe=yellow!50!black]
\textbf{Idée} $\rightarrow$ \textbf{pseudo-code} $\rightarrow$ \textbf{code étape par étape} $\rightarrow$ \textbf{tests} $\rightarrow$ \textbf{corrections}
\end{tcolorbox}

\begin{tabularx}{\linewidth}{>{\bfseries}p{3.2cm}Y Y}
\toprule
Type d'erreur & Signe & Exemple\\
\midrule
Syntaxe & Message avant l'exécution & \texttt{if x > 3} (deux-points oublié)\\
Exécution & Programme qui s'arrête (Traceback) & \texttt{10 / 0}\\
Logique & Résultat faux sans message & \texttt{moyenne = a + b / 2}\\
\bottomrule
\end{tabularx}

\begin{python}[Modèle de calculatrice]
while True:
    print("1 : addition   2 : soustraction   0 : quitter")
    choix = input("Choix : ")
    if choix == "0":
        break
    a = float(input("Premier nombre : "))
    b = float(input("Second nombre : "))
    if choix == "1":
        print("Resultat :", a + b)
    elif choix == "2":
        print("Resultat :", a - b)
\end{python}

%==================================================================
\newpage
\section*{Annexe B : Fiche de pratique guidée et grille du projet}
\entete{green}{Pratique guidée : mon programme}

\subsection*{3a. Pseudo-code de mon projet}
Mon projet : \cb Calculatrice \quad \cb Jeu de devinette\\
Étapes (à ranger dans l'ordre) :\\
\lig \lig \lig \lig

\subsection*{3b. Jeu d'essais}
\begin{tabularx}{\linewidth}{>{\bfseries}p{3cm}Y Y Y}
\toprule
Entrée & Résultat attendu & Résultat obtenu & Conforme ?\\
\midrule
 & & & \\ & & & \\ & & & \\
\bottomrule
\end{tabularx}

\subsection*{Grille d'évaluation du projet}
\begin{tabularx}{\linewidth}{>{\bfseries}p{3.4cm}Y p{2cm}}
\toprule
Critère & Description & Points\\
\midrule
Fonctionnement & Le programme fait ce qui est demandé & /5\\
Structure & Variables, condition et boucle utilisées à bon escient & /5\\
Tests & Jeu d'essais complet, erreurs corrigées & /5\\
Lisibilité & Commentaires, noms clairs, indentation & /5\\
\bottomrule
\end{tabularx}

\begin{tcolorbox}[colback=yellow!10,colframe=yellow!50!black,title={Cartes d'aide à utiliser dans l'ordre},fonttitle=\bfseries\small]
\textbf{Aide 1 :} regarde le pseudo-code du 3a et les modèles des séances précédentes.\\
\textbf{Aide 2 :} Quelle est la prochaine étape du pseudo-code ? Quelle structure faut-il ?\\
\textbf{Aide 3 :} demande à l'enseignant d'écrire la première ligne avec toi.
\end{tcolorbox}

%==================================================================
\newpage
\section*{Annexe C : Niveau A (en dessous du niveau)}
\entete{orange}{Niveau A : calculatrice à deux opérations}

\textbf{Banque de mots :} input, float, if, elif, print, +, -.

\begin{python}[niveauA.py]
a = _____(input("Premier nombre : "))
b = _____(input("Second nombre : "))
op = input("Operation (+ ou -) : ")
_____ op == "+":
    print("Resultat :", a _____ b)
_____ op == "-":
    print("Resultat :", a _____ b)
\end{python}

\subsection*{Activité : trois calculs}
\begin{tabularx}{\linewidth}{>{\bfseries}p{4cm}Y Y}
\toprule
Calcul & Résultat attendu & Résultat obtenu\\
\midrule
$5 + 3$ & & \\ $10 - 4$ & & \\ $2{,}5 + 1{,}5$ & & \\
\bottomrule
\end{tabularx}

%==================================================================
\newpage
\section*{Annexe D : Niveau B (au niveau)}
\entete{blue}{Niveau B : calculatrice complète ou jeu}

\textbf{Option 1 : calculatrice.} Menu en boucle avec les quatre opérations, option quitter, division par zéro refusée avec un message.\\
\textbf{Option 2 : jeu de devinette.} Nombre secret, indices «~plus grand~» ou «~plus petit~», compteur d'essais, rejouer.

\begin{tabularx}{\linewidth}{>{\bfseries}p{3.4cm}Y Y Y}
\toprule
Cas testé & Résultat attendu & Obtenu & Corrigé ?\\
\midrule
Cas normal & & & \\
Cas limite (0, négatif) & & & \\
Saisie absurde (lettre) & & & \\
\bottomrule
\end{tabularx}

\subsection*{Bilan des erreurs rencontrées}
\begin{tabularx}{\linewidth}{>{\bfseries}p{2.8cm}Y Y}
\toprule
Type d'erreur & Où ? & Comment corrigée ?\\
\midrule
 & & \\ & & \\
\bottomrule
\end{tabularx}

%==================================================================
\newpage
\section*{Annexe E : Niveau C (au-dessus du niveau)}
\entete{violet}{Niveau C : projet structuré}

\begin{python}[Aide : fonction et gestion d'erreur]
def addition(a, b):
    return a + b

try:
    a = float(input("Nombre : "))
except ValueError:
    print("Ce n'est pas un nombre.")
\end{python}

\begin{enumerate}
 \item Ajoute une fonction par opération et utilise-les dans la boucle.
 \item Garde un historique des calculs dans une liste (\texttt{historique.append(...)}) et affiche-le à la fin.
 \item Explique pourquoi une fonction améliore la lisibilité : \dotfill
\end{enumerate}

\textbf{Défi bonus :} ajoute les puissances et le reste de la division (\texttt{**} et \texttt{\%}).

%==================================================================
\newpage
\section*{Annexe F : Corrections et attendus (enseignant)}

\subsection*{Annexe B}
Pseudo-code attendu : boucle, menu, saisie de deux nombres, calcul selon l'opération, affichage, quitter. Barème : voir la grille.

\subsection*{Annexe C}
\texttt{float}, \texttt{float}, \texttt{if}, \texttt{+}, \texttt{elif}, \texttt{-}. Résultats : 8, 6, 4.

\subsection*{Annexe D}
Division par zéro : \texttt{if b == 0: print("Division impossible")}. Cas limites : 0, négatifs. Saisie absurde : \texttt{ValueError} (à gérer au niveau C). Les trois types d'erreurs doivent être nommés correctement.

\subsection*{Annexe E}
Structure attendue : fonctions \texttt{addition}, \texttt{soustraction}, \texttt{multiplication}, \texttt{division} ; liste \texttt{historique} ; \texttt{try / except ValueError}.\\
\textbf{Critères de réussite :} programme qui ne plante pas sur une saisie absurde, fonctions utilisées, historique correct.

%==================================================================
\newpage
\section*{Annexe G : Exit ticket et grille d'observation}
\entete{gray}{Exit ticket : programme complet}

\begin{enumerate}
 \item Un programme s'exécute mais affiche un mauvais résultat sans message d'erreur : de quel type d'erreur s'agit-il ?\\ \lig
 \item Cite deux cas que tu testerais pour vérifier une calculatrice.\\ \lig
 \item \autoeval
\end{enumerate}
\textbf{Attendus :} 1. Une erreur de logique. 2. Par exemple une division par zéro, un nombre négatif, une saisie qui n'est pas un nombre.

\bigskip
\grille{Pseudo-code}{Jeu d'essais}{Erreurs corrigées}

\end{document}
